% !TeX program = pdflatex
% Beispiel für die Befehle von schreibvorlage.sty: eine Seite mit großer
% Lineatur, eine mit kleiner und eine ohne Lineatur.
\documentclass{article}
\usepackage[gross]{schreibvorlage}   % oder [klein] oder [keine]

\begin{document}
\anweisung{Spure nach und schreibe ab:}
\reihe{Aa}                           % Aa schwarz, dann grau bis zum Zeilenende
\zeile{}                             % leere Zeile
\reihe{Mm}
\zeile{}
\anweisung{Schreibe die Wörter ab:}
\zeile{\vor{Mama} \nach{Mama Mama}}  % frei gemischt
\zeile{}
\zeile{\vor{Lama} \nach{Lama Lama}}
\zeile{}
\zeile{\vor{Alma} \nach{Alma Alma}}
\zeile{}
\zeile{}

\lineatur{klein}                     % ab einer neuen Seite kleine Lineatur
\anweisung[2]{In der kleinen Lineatur. Eine längere Anweisung belegt
  mehrere Zeilenplätze, damit die Seite im Raster bleibt.}
\reihe{Ll}
\zeile{}
\reihe{Lamm}
\zeile{}
\foreach \i in {1,...,11} {\zeile{}}

\lineatur{keine}                     % Seiten mit gewöhnlichem Text
\section*{Ohne Lineatur}
Gewöhnlicher Text, und mit {\schulschrift Mama} auch in Schulschrift.
\end{document}
